% ============================================================================
% Active-Ledger — FYP Presentation Study Guide (LaTeX)
% Purpose: Help prepare for the final jury / viva by turning the slide deck
% into a full speaking script + deep "why/how" notes + Q&A bank.
%
% Source slides: friends experimentation/Active_Ledger_FYP_Jury (2).html
% Paper reference: docs/paper-from-gemini.tex
% ============================================================================
\documentclass[11pt,a4paper]{article}

\usepackage[T1]{fontenc}
\usepackage[margin=1in]{geometry}
\usepackage{amsmath,amssymb}
\usepackage{graphicx}
\usepackage{booktabs}
\usepackage{enumitem}
\usepackage{hyperref}
\usepackage{xspace}

\setlist[itemize]{leftmargin=*, itemsep=2pt, topsep=2pt}
\setlist[enumerate]{leftmargin=*, itemsep=2pt, topsep=2pt}

% --- Key acronyms (for consistency with paper/slides) ------------------------
\newcommand{\actledger}{\textsc{Active-Ledger}\xspace}
\newcommand{\poc}{\textsc{PoC}\xspace}
\newcommand{\fedavg}{\textsc{FedAvg}\xspace}
\newcommand{\krum}{\textsc{Krum}\xspace}
\newcommand{\mkrum}{\textsc{Multi-Krum}\xspace}
\newcommand{\median}{\textsc{Median}\xspace}
\newcommand{\tmean}{\textsc{Trimmed-Mean}\xspace}
\newcommand{\bulyan}{\textsc{Bulyan}\xspace}

\newcommand{\apbf}{APB-F1\xspace}

\title{\actledger: Full FYP Presentation Preparation Guide\\
\large From basics to ``inside-out'' understanding}
\author{(Personal study notes)}
\date{April 2026}

\begin{document}
\maketitle

\tableofcontents
\newpage

% ============================================================================
\section{How to use this guide (quick)}

\begin{itemize}
\item Treat each section as \textbf{speaker notes for one slide}. The structure mirrors the HTML deck.
\item For each slide:
  \begin{itemize}
  \item Read the \textbf{One-line goal}.
  \item Memorise the \textbf{Core story} in plain language.
  \item Understand the \textbf{Why it is designed this way} bullets.
  \item Practice the \textbf{Likely questions} at the end.
  \end{itemize}
\item The final sections contain a \textbf{deep Q\&A bank}. This is where viva marks are won.
\end{itemize}

\section{Project in one paragraph (opening script)}

\textbf{Script (30--45 seconds).}
Clinical ECG data is sensitive and cannot be pooled across hospitals, so federated learning is attractive. But real hospitals have non-independent and identically distributed (non-IID) data, extreme class imbalance, and a realistic risk of Byzantine participants that poison updates or inject backdoors. \actledger is a blockchain-orchestrated FL framework that defends the minority class that matters clinically---atrial premature beats (APB)---by using a behavioural trust score, Proof-of-Contribution (PoC), computed from on-chain accuracy history. The same PoC also gates synthetic augmentation requests to close a separate attack surface: poisoning through the data-generation pipeline. The ledger provides tamper-evident receipts for both aggregation and synthetic generation so decisions are auditable.

\section{What is being claimed (what to repeat during questions)}

\begin{itemize}
\item \textbf{Primary claim (quantitative).} Across three attacks, \poc is the only defence that keeps \apbf above the clinical threshold of $0.75$.
\item \textbf{Secondary claim (system/security).} Synthetic augmentation is an attack surface; trust-gating plus on-chain receipts closes it in a way robust aggregation alone cannot.
\item \textbf{Blockchain claim (bounded).} The chain is not ``doing learning'' and PoC is not a cryptographic proof; the ledger provides \emph{tamper-evident history} and \emph{auditable approvals} that enable accountable governance.
\end{itemize}

% ============================================================================
\section{Slide-by-slide preparation (mirrors the HTML deck)}

\subsection{Slide 0 --- Title / framing}
\textbf{One-line goal.} State the problem domain and the three pillars: FL, governance, augmentation.

\textbf{Core story (simple).}
This project builds an FL system for ECG classification that is robust to adversarial clients and also provides an auditable governance trail using blockchain.

\textbf{Key pillars to name.}
\begin{itemize}
\item \textbf{Federated learning:} data stays at hospitals.
\item \textbf{Behavioural trust score (PoC):} uses multi-round history rather than single-round geometry.
\item \textbf{Trust-gated synthetic augmentation:} diffusion generator is helpful but unsafe unless governed.
\end{itemize}

\textbf{Likely question.} ``What is new here compared to standard robust aggregation?''
\begin{itemize}
\item Robust aggregation filters model updates; it does not govern \emph{synthetic data requests}.
\item PoC is history-based (behavioural) rather than geometry-based (single round).
\end{itemize}

\subsection{Slide 1 --- Table of contents}
\textbf{One-line goal.} Show the jury the roadmap so the talk feels controlled.

\textbf{How to speak it.} ``First the clinical problem and research gap, then the architecture, then PoC and diffusion, then experiments and results, and finally timeline and conclusion.''

\subsection{Slide 2 --- Problem statement}
\textbf{One-line goal.} Explain the \emph{three simultaneous failures} that make clinical FL hard.

\textbf{Core story.}
\begin{enumerate}
\item \textbf{Non-IID hospital data:} sites have different demographics/equipment; honest updates look like outliers.
\item \textbf{Class imbalance:} rare arrhythmias (APB) are clinically important but easy to erase.
\item \textbf{Governance gap:} typical FL assumes a trusted server; without receipts, post-mortem accountability is weak.
\end{enumerate}

\textbf{Why this matters clinically.}
\begin{itemize}
\item Missing APB is dangerous (delays atrial fibrillation workup).
\item False APB causes unnecessary consults, so precision matters as well as recall.
\end{itemize}

\textbf{Likely question.} ``Why focus on APB rather than overall accuracy?''
\begin{itemize}
\item Because overall metrics can stay high even when minority class fails.
\item Clinical deployability depends on not missing the dangerous minority class.
\end{itemize}

\subsection{Slide 3 --- Prior work}
\textbf{One-line goal.} Show that existing pieces exist, but none solves the combined problem.

\textbf{How to explain quickly.}
\begin{itemize}
\item \fedavg: no defence.
\item \krum/\mkrum/\median/\tmean/\bulyan: geometry-only, per-round; harmed by non-IID and bypassed by sleeper/backdoor behaviour.
\item BlockFL / blockchain FL: often passive logging or coordination, not active defence and not augmentation governance.
\item Time-series GANs: can generate signals, but governance of who can request/use synthetic samples is missing.
\end{itemize}

\textbf{Likely question.} ``Why not just use reputation FL (Kang et al.)?''
\begin{itemize}
\item Reputation FL focuses on identity and incentives; this work focuses on a minimal behavioural score and the augmentation attack surface.
\item The key novelty is \emph{closing the augmentation stealth channel} with verifiable approvals.
\end{itemize}

\subsection{Slide 4 --- Research gap matrix}
\textbf{One-line goal.} Make the gap legible as a checklist of capabilities.

\textbf{How to defend the claim.}
The gap is not ``blockchain'' alone or ``robust aggregation'' alone; it is the combination:
\begin{itemize}
\item provenance + robust aggregation + behavioural memory + governed augmentation + closing the augmentation stealth channel + rare-class threshold across attacks.
\end{itemize}

\textbf{Important caution.}
Avoid over-claiming. The paper intentionally frames blockchain as a tamper-evident log and approval ledger, not a new consensus protocol.

\subsection{Slide 5 --- System architecture}
\textbf{One-line goal.} Walk through the system with clear data flow and decision points.

\textbf{Architecture in three layers.}
\begin{itemize}
\item \textbf{Clinical edge (clients).}
  \begin{itemize}
  \item 10 clients (8 honest, 2 Byzantine), Dirichlet $\alpha=0.5$ partitioning, 40 rounds, 3 local epochs.
  \item CNN--LSTM classifier on 360-sample windows at 360Hz.
  \item Private data stays local; only weights and metrics are shared.
  \end{itemize}
\item \textbf{Off-chain orchestrator (server).}
  \begin{itemize}
  \item Receives updates; logs receipts on-chain.
  \item Computes PoC score using EMA of \emph{validation accuracy history} and participation.
  \item Aggregates \textbf{top-$K=7$} clients by PoC (weighted by sample size).
  \item Approves synthetic requests only if PoC $\ge \tau$.
  \end{itemize}
\item \textbf{Ledger (Ganache Ethereum + Solidity contract).}
  \begin{itemize}
  \item Stores events: update logs, round completion, synthetic requested/approved/generated.
  \item Stateless/event-only design: emphasis on receipts and auditability (not training).
  \end{itemize}
\end{itemize}

\textbf{Why Ganache? (important viva topic).}
\begin{itemize}
\item It provides a fully controlled local Ethereum environment for a research prototype.
\item Deterministic, fast, no real cost, no external dependencies during experiments.
\item Allows the same EVM execution model and event logs as mainnet without production overhead.
\end{itemize}

\subsection{Slide 6 --- Proof-of-Contribution (PoC)}
\textbf{One-line goal.} Explain why behaviour across rounds beats geometry in a sleeper/backdoor setting.

\textbf{Definition (plain language).}
PoC is a trust score per client that increases when a client repeatedly reports good validation accuracy and participates consistently, and drops quickly when its behaviour worsens.

\textbf{Equation-level understanding.}
\[
\mathrm{PoC}(i,r)=\mathrm{EMA}_\gamma(\text{accuracy history})\times\text{participation fraction},\quad \gamma=0.7.
\]

\textbf{Why this works against sleeper.}
\begin{itemize}
\item A sleeper is \emph{not} anomalous early; geometry-only filters have no signal.
\item After activation, label flipping hurts the attacker’s own validation accuracy on honest validation data.
\item EMA with $\gamma=0.7$ adapts in a few rounds, pushing attackers out of top-$K$.
\end{itemize}

\textbf{Likely questions.}
\begin{itemize}
\item ``Is PoC a cryptographic proof?'' No. It is a score, not a consensus protocol.
\item ``Could an attacker fake the accuracy?'' In this design, accuracy is evaluated against a public proxy subset held by the orchestrator (paper detail); a production design would require stronger attestation or secure evaluation.
\item ``Why top-$K=7$?'' It trades robustness and stability: enough diversity to avoid overfitting to a few sites, but excludes low-trust clients.
\end{itemize}

\subsection{Slide 7 --- Trust-gated latent diffusion}
\textbf{One-line goal.} Introduce the augmentation stealth channel and explain why it is upstream of aggregation.

\textbf{The attack (augmentation stealth channel).}
\begin{itemize}
\item Byzantine client requests synthetic data but flips the requested label.
\item The generator produces samples from the wrong class manifold.
\item Those samples are injected into local training \emph{before} aggregation; robust aggregators cannot ``undo'' poisoned data.
\end{itemize}

\textbf{The defence (trust gate + receipts).}
\begin{itemize}
\item Synthetic request is approved only if PoC $\ge \tau=0.4$.
\item After sleeper activation, attacker’s PoC drops; requests are rejected.
\item The ledger records \texttt{Requested/Approved/Generated} so every sample batch is auditable.
\end{itemize}

\textbf{Why 1D diffusion for ECG?}
\begin{itemize}
\item ECG is a 1D time-series; 1D UNet fits the structure and is efficient.
\item Diffusion is empirically more stable than GANs for generation (less mode collapse, better diversity).
\item Conditioning on class label supports targeted rare-class augmentation.
\end{itemize}

\textbf{Likely question.} ``Why not use GAN/TimeGAN like the proposal?''
\begin{itemize}
\item The project evolved: diffusion was used because it yielded higher quality and more stable training on the target setup.
\item In the final story, the key contribution is governance and security of augmentation, not claiming diffusion is universally best.
\end{itemize}

\subsection{Slide 8 --- Experimental setup}
\textbf{One-line goal.} Make the evaluation replicable and defensible.

\textbf{Dataset and preprocessing.}
\begin{itemize}
\item MIT-BIH Arrhythmia Database, AAMI EC57 mapping.
\item 360-sample windows at 360Hz, standard normalisation.
\item Five classes: Normal, LBBB, RBBB, APB, PVC.
\end{itemize}

\textbf{Federation parameters.}
\begin{itemize}
\item 10 clients; 2 Byzantine (20\%).
\item 40 rounds; 3 local epochs; Adam LR $10^{-3}$; batch size 128.
\item Dirichlet concentration $\alpha=0.5$ (non-IID stress test).
\item PoC: EMA $\gamma=0.7$; trust threshold $\tau=0.4$; top-$K=7$.
\end{itemize}

\textbf{Attacks (know them precisely).}
\begin{itemize}
\item \textbf{Gaussian poisoning:} weights replaced by uniform noise each round.
\item \textbf{Static label-flip:} targeted pairwise swap Normal $\leftrightarrow$ LBBB.
\item \textbf{Sleeper:} behaves honestly until round 15, then backdoors APB $\rightarrow$ LBBB with backdoor-critical-layer fine-tuning.
\end{itemize}

\subsection{Slide 9 --- Results}
\textbf{One-line goal.} Make the jury care about the APB safety story.

\textbf{The table to memorise (APB-F1).}
\begin{itemize}
\item \fedavg: collapses under Gaussian (APB-F1 = 0.000).
\item \krum: collapses under sleeper (APB-F1 \(\approx 0.306\)).
\item \poc: only method above 0.75 across all three attacks (worst \(\approx 0.786\)).
\item Trust-gated diffusion: APB-F1 \(\approx 0.879\) vs blind baseline \(\approx 0.857\).
\end{itemize}

\textbf{How to answer ``but FedAvg looks good on macro-F1''.}
\begin{itemize}
\item Macro-F1 can hide minority-class failures.
\item In non-IID + low attacker ratio, mean aggregation may dilute poison on common classes.
\item APB is the deployability bottleneck; therefore \apbf is the main metric.
\end{itemize}

\textbf{Gas-cost slide caveat.}
If asked: the paper frames gas accounting as out-of-scope for the main claim; the prototype uses event logs rather than storage-heavy writes. For production, a permissioned chain or L2 could be considered.

\subsection{Slide 10 --- Timeline}
\textbf{One-line goal.} Demonstrate research maturity and iteration.

\textbf{What to highlight.}
The project iterated from baseline FL $\rightarrow$ blockchain integration $\rightarrow$ PoC $\rightarrow$ diffusion gate $\rightarrow$ formal threat model and stronger baselines after ACISP critique.

\subsection{Slide 11 --- Team and work division}
\textbf{One-line goal.} Show clear ownership and collaboration.

\textbf{Prepare one sentence each about contributions.}
\begin{itemize}
\item Blockchain and orchestration: smart contract, event schema, Web3 integration.
\item ML and research: PoC, threat model, attacks, diffusion integration, paper.
\item Data and evaluation: preprocessing, partitioning, metrics, plots, consistency.
\end{itemize}

\subsection{Slide 12 --- ACISP critique and ESORICS response}
\textbf{One-line goal.} Show ability to take reviewer feedback and improve scientifically.

\textbf{Key lesson.}
Most critique categories are ``research hygiene'': baselines, threat model clarity, symbol hygiene, consistency between prose and figures, and bounded claims about blockchain.

\subsection{Slide 13 --- Proposal vs final}
\textbf{One-line goal.} Explain evolution without sounding like scope drift.

\textbf{Explain as ``research refinement''.}
\begin{itemize}
\item Some proposed items were replaced with stronger or more relevant ones (GAN $\rightarrow$ diffusion, 3 clients $\rightarrow$ 10 clients, weak baselines $\rightarrow$ robust baselines).
\item New insight emerged: the augmentation stealth channel was not obvious at proposal time.
\end{itemize}

\subsection{Slide 14 --- Conclusions}
\textbf{One-line goal.} End with three clean takeaways and believable future work.

\textbf{The three findings.}
\begin{itemize}
\item Behavioural memory (PoC) beats geometry for sleeper/backdoor behaviour.
\item Augmentation governance is necessary; trust-gated diffusion improves safety and lowers backdoor success.
\item Blockchain earns its place via tamper-evident receipts and governed approvals (bounded claim).
\end{itemize}

\subsection{Slide 15 --- Thank you / final bullets}
\textbf{One-line goal.} Leave a ``numbers + story'' imprint: 0.75 threshold, 0.879, 0.857, 0.000, 0.306.

% ============================================================================
\newpage
\section{Deep understanding: key design decisions (the ``why'' section)}

\subsection{Why CNN--LSTM for ECG classification?}
\begin{itemize}
\item \textbf{Signal structure.} ECG has local morphology (QRS complexes) and temporal dependencies (beat-to-beat context). CNN captures local patterns; LSTM captures sequence context.
\item \textbf{Practicality.} The architecture is compact, trains reliably, and is light enough for FL clients.
\item \textbf{Why not Transformers?} Transformers are heavier, need more data, and introduce more hyperparameter sensitivity; the project targeted a stable baseline to test security mechanisms.
\item \textbf{Why not pure CNN?} Pure CNN can work, but LSTM adds temporal robustness in noisy clinical segments.
\end{itemize}

\subsection{Why 1D latent diffusion for ECG synthesis?}
\begin{itemize}
\item \textbf{Diffusion stability.} Diffusion training tends to be more stable than GANs (less collapse), important for research timelines.
\item \textbf{1D fit.} ECG is 1D; a 1D UNet is efficient and matches inductive bias.
\item \textbf{Class-conditional generation.} Conditioning allows targeted rare-class augmentation.
\item \textbf{Why latent diffusion (vs pixel-space diffusion).} Latent versions reduce compute; in a research prototype, this matters.
\end{itemize}

\subsection{Why reuse the same model across experiments?}
\begin{itemize}
\item \textbf{Controlled comparison.} If the classifier changes across defences/attacks, it becomes unclear whether differences are due to security mechanism or model capacity.
\item \textbf{Fairness.} All defences are tested on identical learning dynamics.
\item \textbf{Research goal.} The contribution is a security/governance mechanism; the base classifier should be stable, not ``cherry-picked''.
\end{itemize}

\subsection{Why evaluate 7 aggregation defences?}
\begin{itemize}
\item \textbf{Community expectation.} Robust FL is judged against these canonical baselines.
\item \textbf{Coverage.} Geometry-only filters (Krum family, coordinate-wise) and composed defences (Bulyan) represent common families.
\item \textbf{Narrative clarity.} It shows the difference between geometry-only vs behavioural history.
\end{itemize}

\subsection{Why Ganache and an event-only smart contract?}
\begin{itemize}
\item \textbf{Prototype speed and determinism.} Ganache runs locally, enabling repeatable experiments.
\item \textbf{Receipts vs state.} Events are lightweight and preserve an append-only audit trail without expensive persistent storage writes.
\item \textbf{Bounded claim.} The ledger is used for auditability and governance, not to claim decentralised training.
\end{itemize}

\subsection{Why not just use a signed log instead of blockchain?}
\begin{itemize}
\item A signed log still requires a trust anchor for who maintains it and can be disputed in multi-party settings.
\item Blockchain provides a shared, tamper-evident timeline and public verifiability within the consortium model.
\item For a research prototype, it demonstrates accountable governance; production could also use permissioned ledgers.
\end{itemize}

% ============================================================================
\newpage
\section{Viva / Jury Q\&A bank (practice answers)}

\subsection{High-probability questions (top 20)}
\begin{enumerate}
\item \textbf{What is the exact research contribution?}\\
Rare-class robustness under attack (APB-F1 across three attacks) + governed augmentation (trust-gated diffusion) + auditable receipts enabling accountable decisions.

\item \textbf{Why is APB the primary metric?}\\
APB is rare and clinically important; macro metrics can hide minority failure; a clinically deployable system cannot sacrifice the rare dangerous class.

\item \textbf{What does PoC measure, exactly?}\\
Historical validation accuracy (EMA) multiplied by participation; it is a behavioural reliability score, not a cryptographic proof.

\item \textbf{How does PoC stop sleeper attacks?}\\
Sleeper behaves honestly early; after activation, its own validation accuracy drops; EMA reacts and reduces score; attackers fall out of top-$K$.

\item \textbf{Why can’t Krum stop the sleeper attack?}\\
Krum is single-round geometry; sleeper updates are crafted to look geometrically normal; there is no early anomaly, and post-activation updates can remain close to honest ones in weight space.

\item \textbf{What is the augmentation stealth channel?}\\
Poisoning via synthetic data requests: attacker requests minority samples but flips labels; poison enters training data before aggregation.

\item \textbf{Why is the blockchain needed?}\\
To store tamper-evident receipts for updates and synthetic approvals; enables accountability and prevents history rewriting in consortium setting.

\item \textbf{Why Ganache instead of public Ethereum?}\\
Local deterministic research environment; avoids real cost and latency; still uses EVM semantics and event logs.

\item \textbf{Why CNN--LSTM?}\\
CNN captures local ECG morphology; LSTM captures temporal dependencies; compact and reliable for FL.

\item \textbf{Why diffusion, why 1D UNet?}\\
Stable generation, good diversity; ECG is 1D; UNet is standard for diffusion and efficient in 1D.

\item \textbf{Why were all experiments done on the same classifier?}\\
To isolate the effect of the defence; changing classifier confounds results.

\item \textbf{What are the three attacks precisely?}\\
Gaussian noise weights; Normal$\leftrightarrow$LBBB static label swap; sleeper APB$\rightarrow$LBBB backdoor via critical-layer flip after round 15.

\item \textbf{Why choose Dirichlet $\alpha=0.5$?}\\
Common non-IID stress setting; creates strong heterogeneity across clients.

\item \textbf{What is the trust threshold $\tau=0.4$?}\\
Minimum PoC required to approve synthetic requests; chosen to block post-activation attackers once their score drops.

\item \textbf{Why top-$K=7$?}\\
Exclude low-trust clients while keeping enough diversity; balances robustness and stability.

\item \textbf{What is BSR and why is it needed?}\\
Backdoor Success Rate captures residual backdoor capacity even when APB-F1 looks healthy; security metric for augmented setting.

\item \textbf{What are the limitations?}\\
Single seed per cell; fixed 20\% attacker ratio; limited attack coverage; gas cost analysis not the main contribution; no cross-dataset generalisation.

\item \textbf{What would be done differently for production?}\\
Permissioned ledger/L2; stronger attestation for metrics; more seeds; higher attacker ratios; adaptive adversary evaluation; broader datasets.

\item \textbf{Does the blockchain slow training?}\\
In the prototype overhead is small (event logging); latency is not claimed as a contribution.

\item \textbf{If a client lies about accuracy, does PoC fail?}\\
If accuracy were self-reported, yes; the paper clarifies accuracy is evaluated against an orchestrator-held public proxy subset. Production would require stronger evaluation/attestation.
\end{enumerate}

\subsection{Common follow-ups (short answers)}
\begin{itemize}
\item \textbf{Why not FedProx / personalisation?} The final framing is security/robustness; PoC handles heterogeneity; personalisation was de-scoped.
\item \textbf{Why only 10 clients?} Research prototype + controlled evaluation; enough to show non-IID and 20\% Byzantines while keeping experiments feasible.
\item \textbf{Why only 2 Byzantines?} To match a realistic minority threat model and align with robust aggregation assumptions; future work should test higher ratios.
\end{itemize}

% ============================================================================
\section{Cheat sheets (for last-minute revision)}

\subsection{Numbers to memorise}
\begin{itemize}
\item Clinical threshold: \textbf{APB-F1 = 0.75}.
\item Worst-case \poc across attacks: \textbf{0.786}.
\item \fedavg under Gaussian: \textbf{0.000}.
\item \krum under sleeper: \textbf{0.306}.
\item Diffusion: trust-gated \textbf{0.879} vs blind baseline \textbf{0.857}.
\item Key hyperparameters: 40 rounds, 3 local epochs, batch 128, Adam $10^{-3}$, Dirichlet $\alpha=0.5$, EMA $\gamma=0.7$, trust threshold $\tau=0.4$, top-$K=7$.
\end{itemize}

\subsection{Architecture one-liner}
\textbf{Uploads and metrics are logged as receipts; PoC is computed from accuracy history; top-$K$ aggregation uses PoC; the same PoC gates synthetic requests; the ledger makes both aggregation and augmentation auditable.}

% ============================================================================
% ============================================================================
\newpage
\section{MASTER SECTION --- Complete project, from basics to hero}
\label{sec:master}

\textit{This is the deep-dive section. Read it once end-to-end, slowly.
After this, you should be able to explain any part of the project
without referring to the paper.}

% ----------------------------------------------------------------------------
\subsection{Part A --- The problem in plain English}

\subsubsection{What is the clinical problem?}
A cardiologist wants a reliable model that can look at a short ECG
strip (a small window of a patient's heartbeat) and say which
category the beat belongs to: Normal, Left Bundle Branch Block
(LBBB), Right Bundle Branch Block (RBBB), Atrial Premature Beat
(APB), or Premature Ventricular Contraction (PVC).

Two things are uncomfortable:
\begin{enumerate}
\item Patient data cannot leave the hospital (privacy laws, HIPAA-
style regulations). That rules out pooling all data in one place.
\item Rare but dangerous heartbeats (especially APB) are the
clinically important ones. If the model gets APB wrong, a doctor
either misses an early atrial fibrillation signal, or does an
unnecessary consult.
\end{enumerate}

So the model has to be trained \textbf{without centralising data}
and also has to \textbf{not fail on rare classes}.

\subsubsection{Why not just train a classical model in one hospital?}
\begin{itemize}
\item A single hospital's data is not diverse enough (age, equipment, ethnicity, protocols vary across sites).
\item The rare classes in one hospital might be extremely rare, so a local model overfits.
\item To generalise across the real world, the model has to see many distributions.
\end{itemize}

\subsubsection{Why does this become a security problem?}
Once multiple hospitals share a model-training system:
\begin{itemize}
\item Any compromised hospital can upload corrupted weight updates. That is a \textbf{Byzantine} client.
\item The attacker can degrade overall accuracy or, more dangerously, silently make the model wrong on specific classes (a \textbf{backdoor}).
\item In a hospital setting, a silent backdoor on APB is not a ``bug'', it is a patient safety risk.
\end{itemize}

That is why the project is not ``just FL'' but ``FL + defence + auditability''.

% ----------------------------------------------------------------------------
\subsection{Part B --- Full system flow (end-to-end trace of one round)}

Imagine a single federated round, $r$. Nothing is abstract here; this is exactly what the system does:

\begin{enumerate}
\item \textbf{Selection.}
The server looks at the PoC score $\mathrm{PoC}(i, r)$ for every
client $i$, ranks all clients, and picks the top $K=7$. The other
three are skipped for this round. For round $r=0$ the PoC score is
initialised at $0$ for everyone, so in the very first rounds
selection behaves almost like uniform selection.

\item \textbf{Local training on each selected client.}
\begin{itemize}
\item The client receives the current global model weights $\theta_r$.
\item It runs \textbf{3 local epochs} of Adam (LR $10^{-3}$, batch $128$) on its own training split.
\item If the client has rare-class deficit (fewer than $\sim 50$ samples of APB, say), it files a \texttt{SyntheticRequested} event on the ledger asking for synthetic help for that class.
\end{itemize}

\item \textbf{Synthetic request handling.}
The approval daemon reads the request on-chain:
\begin{itemize}
\item If the requester's current PoC is $\ge \tau = 0.4$: approve.
\item Else: reject (recorded on-chain anyway).
\end{itemize}
If approved, the 1D latent diffusion generator runs 30 DDPM steps with the requested class label as the conditioning signal, returns $\sim 500$ synthetic ECG windows, and emits \texttt{SyntheticGenerated}. The client then retrains using both real and synthetic samples for that round.

\item \textbf{Upload and receipt.}
Each selected client uploads its new weights $\theta_i^{r+1}$ to the server. The server writes an \texttt{UpdateLogged} event on-chain with the client ID, round, the SHA-256 hash of the weights, and the client's self-reported training-round accuracy.

\item \textbf{Orchestrator scoring.}
The orchestrator holds a small, public, anonymised \emph{proxy subset} of MIT-BIH. It evaluates each submitted model on that proxy and gets a validation accuracy $\mathrm{acc}_{i,r}$ per client. This accuracy is used to update the PoC EMA:
\[
\mathrm{EMA}_{\gamma}\big(\mathrm{acc}_{i,:r}\big) \gets \gamma \cdot \mathrm{acc}_{i,r} + (1-\gamma) \cdot \mathrm{EMA}_{\gamma}\big(\mathrm{acc}_{i,:r-1}\big),\quad \gamma=0.7.
\]
Multiply by the participation fraction: \emph{(rounds with a valid \texttt{UpdateLogged}) / (total rounds so far)}. Multiply these two numbers to get $\mathrm{PoC}(i, r+1)$. Clamp into $(10^{-6}, 1-10^{-6})$.

\item \textbf{Aggregation.}
The server aggregates the top-$K$ updates weighted by local sample size (honest non-IID clients with more data legitimately count more). The new global model $\theta_{r+1}$ is formed.

\item \textbf{Round completion receipt.}
The server writes a \texttt{RoundCompleted} event on-chain with the hash of $\theta_{r+1}$. This is what makes the history auditable: a third party can later verify which models were accepted on which round.

\item \textbf{Broadcast.}
$\theta_{r+1}$ is sent back to all clients for round $r+2$.
\end{enumerate}

Repeat for 40 rounds. Evaluate final per-class F1 on a held-out 15\% test split per client.

\textit{If you can recite this 8-step loop fluently in the viva, you can defend the architecture.}

% ----------------------------------------------------------------------------
\subsection{Part B$^\star$ --- Federated Learning (conceptual deep dive)}
\label{sec:fl-conceptual}

\textit{This subsection is the ``conceptual FL'' layer: what FL is, why FedAvg works,
why it fails under non-IID, what client drift is, and why geometry-only defences behave the way they do.}

\subsubsection{B$^\star$.1 --- What federated learning is (one clean definition)}
Federated learning (FL) is a training paradigm where multiple clients (hospitals)
train a shared model collaboratively without sharing raw data. Each client runs
local optimisation on its private data and sends only a model update (weights or
gradients) to a server. The server aggregates these updates into a new global
model and broadcasts it back. This repeats across rounds.

\subsubsection{B$^\star$.2 --- Why \fedavg works in the IID world (intuition)}
In the simplest case, every hospital's data distribution is identical
(independent and identically distributed, IID). Then each client's gradient is
an unbiased estimate of the same underlying ``true'' gradient of the global
objective. Averaging many such gradients reduces noise (variance), so the mean
update points in a useful direction. This is why \fedavg is a strong baseline
when the assumptions hold and there are no adversaries.

In plain language: if everyone is learning from the same kind of patients, then
``average of what everyone learned'' is a good approximation of what would have
happened if all data were pooled.

\subsubsection{B$^\star$.3 --- What changes under non-IID data (the core pain)}
In hospitals, data is almost never IID:
\begin{itemize}
\item Site A might have older patients, Site B younger patients.
\item One site might see more APB cases due to referral patterns.
\item Equipment and preprocessing pipelines differ.
\end{itemize}

Under non-IID, each client is optimising a \emph{different local objective}.
Therefore, local gradients can point in different directions. This creates two
important phenomena:
\begin{itemize}
\item \textbf{Legitimate outliers.} Honest clients can look far from the mean because their data is genuinely different.
\item \textbf{Slower or unstable convergence.} Averaging conflicting updates can move the global model in a direction that helps majority sites and hurts minority sites.
\end{itemize}

\subsubsection{B$^\star$.4 --- Client drift (what it is and why it matters)}
Client drift means: after several local steps, a client's model can move toward
its local optimum, which may be far from the global optimum when data is non-IID.
The more local epochs are run, the more drift can happen. Drift is not a security
attack; it is a statistical/optimisation effect. However, drift creates a
security side-effect: it becomes hard to tell ``honest disagreement'' from
``malicious manipulation'' purely from weight geometry.

\subsubsection{B$^\star$.5 --- Why non-IID makes Byzantine defence harder}
Most classical Byzantine-robust FL defences are based on geometry: they try to
remove updates that are ``far'' from the others in the current round.

But in non-IID FL:
\begin{itemize}
\item Some honest clients are far from the mean (because they represent minority demographics or rare cases).
\item Filtering based on distance can accidentally throw away the most clinically valuable clients (the ones who see APB).
\end{itemize}

This is exactly why the project focuses on rare-class preservation: a defence
that protects the average case but discards the rare-class experts is not
clinically deployable.

\subsubsection{B$^\star$.6 --- The difference between geometric and behavioural signals}
\textbf{Geometric signal (single-round).} ``Does this update look like the crowd \emph{right now}?''\\
Used by \krum/\mkrum/\median/\tmean/\bulyan.

\textbf{Behavioural signal (multi-round).} ``Has this client behaved reliably across rounds?''\\
Used by PoC via an EMA of validation accuracy and participation.

The conceptual reason PoC helps is that it introduces the missing dimension:
\textbf{time}. A sleeper attacker is designed to be geometrically normal early,
so single-round geometry has no evidence to act on.

\subsubsection{B$^\star$.7 --- Why sleeper attacks are ``the worst case'' for geometry-only defences}
A sleeper attack behaves honestly for many rounds and activates later.
That means:
\begin{itemize}
\item Before activation, its updates are indistinguishable from honest updates (so geometry-based filters learn nothing).
\item After activation, it can still craft updates to remain close in weight space (especially when only a small subset of backdoor-critical layers is changed).
\end{itemize}

So the defence needs a non-geometric signal. In \actledger that signal is:
``does the client keep achieving good validation accuracy on a clean reference set across time?''

\subsubsection{B$^\star$.8 --- Why the project holds the model constant across experiments}
From a conceptual FL standpoint, a fair defence comparison requires that
\emph{the learning problem be identical} across methods. If the classifier,
generator, optimiser, or data partitioning changes across conditions, then
differences in results can be attributed to those changes rather than the defence.
Holding the model constant makes the experiment about security/governance, not
model capacity.

\subsubsection{B$^\star$.9 --- The shortest ``what FL is'' answer for viva}
\textbf{FL is distributed training where data stays local, clients send model updates, a server aggregates them, and non-IID makes honest clients disagree---which is why robust defence needs either strong statistics or history-based behavioural trust.}

% ----------------------------------------------------------------------------
\subsection{Part B-II --- Non-IID data and the Dirichlet partition (plain English)}
\label{sec:noniid-dirichlet}

\textbf{Simple definition.} ``IID'' means every client's data is drawn from
the same distribution, like cutting one big pizza into equal slices. ``Non-IID''
means each client has a different flavour of data. One hospital might mostly see
Normal beats; another might see many APB cases; another might have older patients
only.

\textbf{Why this matters for FL.} When every client is different, the
``average'' model update doesn't point in the same direction that a central-dataset
model would. Some clients pull the global model toward their majority class;
others pull it toward their niche. The global model becomes a compromise, and
compromises often fail the minority class.

\textbf{Dirichlet partition in one picture.}
Imagine you want to split a dataset among 10 clients and control how unequal the
split is. The Dirichlet distribution with parameter $\alpha$ is the standard
knob:
\begin{itemize}
\item $\alpha = 100$: almost equal slices across clients (close to IID).
\item $\alpha = 0.5$ (this project): strongly skewed, so some clients get a lot
of APB and others get almost none. Realistic and challenging.
\item $\alpha = 0.05$: extreme skew, most clients get only one class.
\end{itemize}
$\alpha$ is effectively a ``fairness knob'' for how balanced the per-class counts
are across clients.

\textbf{Why $\alpha=0.5$ was chosen.} It is the community-standard non-IID stress
test. Strong enough to break naive FedAvg, not so extreme that the experiment
becomes a pathological edge case.

\textbf{One viva-ready line.}
\textit{``Dirichlet $\alpha=0.5$ is our non-IID stress test: each client has a
meaningfully different mix of classes, so honest clients can legitimately disagree,
which is why geometric defences over-penalise minority-class experts.''}

% ----------------------------------------------------------------------------
\subsection{Part B-III --- CNN and LSTM in simple language}
\label{sec:cnn-lstm}

\textbf{ECG as a simple picture.} An ECG window is just a 1D signal: 360 numbers
that go up and down over time (one second at 360\,Hz). The classifier has to
look at this shape and say which class it is.

\textbf{What a CNN does (pattern detector).}
A Convolutional Neural Network (CNN) slides a small ``template'' across the
signal and measures how well the template matches at each position. If a
template looks like a QRS spike, the CNN lights up wherever a QRS spike appears.
Multiple templates are learned automatically; deeper layers combine simple
templates into more complex ones (``a sharp upstroke followed by a wide dip =
premature beat''). \textbf{Analogy:} a CNN is like a moving magnifying glass
that notices small repeating shapes.

\textbf{What an LSTM does (sequence memory).}
A Long Short-Term Memory (LSTM) is a type of recurrent network. It reads the
signal one chunk at a time and maintains a running ``memory state''. This lets
the model decide ``this beat is abnormal because the last two beats were also
abnormal'' rather than judging each beat in isolation. \textbf{Analogy:} the
LSTM is a reader who keeps notes in the margin and uses those notes when deciding
what the next sentence means.

\textbf{Why combine them (CNN-LSTM).}
CNN extracts ``what does this piece of the waveform look like?''
LSTM extracts ``how does the shape evolve across the window?''
Together: local shape + temporal context = more reliable classification,
especially on borderline beats.

\textbf{Why not just Transformer?}
A Transformer is more powerful but hungry for data and compute. In FL, each
client has limited data and compute; Transformers also have many hyperparameters
that would confound a security comparison. CNN-LSTM is the right middle ground.

% ----------------------------------------------------------------------------
\subsection{Part B-IV --- Diffusion models, UNet, DDPM, latent diffusion (plain English)}
\label{sec:diffusion-conceptual}

\textbf{The core idea in one sentence.}
Diffusion models learn to \emph{undo noise}. If you slowly blur a real ECG with
random noise until it becomes garbage, a diffusion model learns the reverse:
it starts from pure noise and step-by-step removes noise until a realistic
sample appears.

\textbf{The forward process (conceptual).}
Think of taking a clean ECG and adding a little noise at each of $T$ steps. By
step $T$ the ECG is indistinguishable from random noise. This is deterministic
and trivial to compute; it doesn't need learning.

\textbf{The reverse process (the model).}
The model is trained to predict, given a noisy signal, how to denoise it by
one step. Apply this $T$ times starting from pure noise and you get a new
synthetic ECG that looks like real data but is not a copy.

\textbf{Why ``DDPM''.}
Denoising Diffusion Probabilistic Models (DDPM) is the specific recipe used
here for how the denoising steps are defined and trained. 30 steps were
enough to recover clinically plausible ECG morphology.

\textbf{What is a UNet?}
A UNet is a neural network shape that first compresses the input, then
expands it back to the same length, with ``skip connections'' between the
compression and expansion sides so detail is preserved. It is the standard
choice for denoisers because the compressed part captures context and the
expansion part restores detail. \textbf{Analogy:} a UNet is like zooming
out to understand a big picture and then zooming back in carefully to draw
the details.

\textbf{Why ``1D''.}
Our signal is a 1D time series, not a 2D image. A 1D UNet slides 1D
convolutions along the time axis. It is lighter and more natural than a 2D
UNet for ECG.

\textbf{Why ``latent'' diffusion.}
Running diffusion directly in signal space is expensive because each of the 30
steps operates on the full 360 numbers. Latent diffusion first compresses the
signal to a shorter representation (``latent space''), runs diffusion there,
and decodes back. Faster training and inference without much quality loss.

\textbf{Why class-conditional.}
We don't want random ECGs; we want ``an APB that looks like an APB''. So we
attach the class label as an extra input channel, and the model learns a
separate denoising trajectory per class.

\textbf{Why diffusion and not GAN?}
GANs train two networks against each other (generator vs discriminator), which
is notoriously unstable (``mode collapse''). Diffusion trains a single denoiser
with a simple objective and is empirically more reliable for medical data.

% ----------------------------------------------------------------------------
\subsection{Part B-V --- Blockchain, smart contracts, Ethereum events, and Ganache (plain English)}
\label{sec:blockchain-conceptual}

\textbf{What a blockchain really is.}
A blockchain is a shared append-only notebook. Many parties hold a copy of it.
New entries are added in ordered ``blocks''. Once an entry is in, rewriting it
would require rewriting every block after it on every copy of the notebook,
which is practically infeasible. That is what ``tamper-evident'' means.

\textbf{What Ethereum adds.}
Ethereum is a blockchain that runs small programs called \emph{smart contracts}.
A contract is just code with variables (state) and functions. When a function
is called, every participant runs the same code and agrees on the result.

\textbf{What a smart contract looks like in this project.}
Ours is deliberately simple. It has no training logic. All it does is emit
\emph{events} when the server reports something:
\begin{itemize}
\item \texttt{UpdateLogged(client, round, weightsHash, reportedAccuracy)}
\item \texttt{RoundCompleted(round, globalModelHash)}
\item \texttt{SyntheticRequested(client, round, requestedClass)}
\item \texttt{SyntheticApproved(client, round, pocAtApproval)}
\item \texttt{SyntheticGenerated(client, round, numSamples, sampleHash)}
\end{itemize}
These events are permanent in the chain history.

\textbf{Why events (cheap) and not storage (expensive)?}
On Ethereum, saving data into contract \emph{state} costs a lot of ``gas''
(the transaction fee). Emitting an \emph{event} is much cheaper. Since we only
need the history to be \emph{readable} later (not queryable from other contracts),
events are the right primitive. That is the 84.6\% gas cost reduction figure.

\textbf{What Ganache is.}
Ganache is a local, private Ethereum simulator. It:
\begin{itemize}
\item Runs entirely on your laptop.
\item Uses the exact same contract execution rules (EVM) as real Ethereum.
\item Is deterministic (you control block timing), fast, and free.
\end{itemize}
\textbf{Analogy:} Ganache is a driving simulator; the car's controls are identical
to a real car, but you can crash without paying damages. Perfect for a research
prototype.

\textbf{Why Ganache and not public Ethereum for a prototype?}
Public Ethereum has real cost, unpredictable latency, and public visibility of
every transaction. For a controlled research experiment, determinism and
speed are essential.

\textbf{Why Ganache and not Hyperledger Fabric?}
Hyperledger is better for a real hospital consortium. But it requires more
setup, its programming model is different from Ethereum, and its documentation
is heavier. For an undergraduate FYP prototype that needs to demonstrate the
concept in a reproducible way, Ethereum + Ganache is the right trade-off. The
paper explicitly notes permissioned chains as a production direction.

\textbf{Key bounded claim.}
The ledger is not doing consensus for training; it is providing an
\emph{auditable trail}. The same system could run on any tamper-evident log
with signed entries. We use Ethereum because its event model is well-understood
and well-tooled.

% ----------------------------------------------------------------------------
\subsection{Part B-VI --- Exponential Moving Average (EMA) in simple language}
\label{sec:ema-conceptual}

\textbf{The idea.} An EMA is a running average that weights recent values more
than older ones. Every new value $x_r$ updates the EMA like this:
\[
\mathrm{EMA}_r = \gamma \cdot x_r + (1 - \gamma) \cdot \mathrm{EMA}_{r-1},
\]
where $\gamma \in (0, 1)$. With $\gamma = 0.7$, the new value contributes 70\%
and the previous running average contributes 30\%.

\textbf{Why not a simple average?}
A simple average over all rounds treats round 1 and round 30 equally. If a client
was honest for rounds 1--10 but malicious from round 11 onward, a simple average
would take forever to catch up. An EMA adapts in a handful of rounds.

\textbf{Why not just ``this round's accuracy''?}
Then one unlucky round could kick out an honest client. An EMA smooths over
noise but still adapts.

\textbf{How this catches a sleeper.}
Before activation, the sleeper's accuracy is high, so its EMA is high. After
activation, its accuracy drops, and within a few rounds its EMA drops enough
that it falls out of the top-$K$ and below the trust threshold.

\textbf{Analogy:} EMA is like a rolling grade that weights the most recent exam
more than the first one.

% ----------------------------------------------------------------------------
\subsection{Part B-VII --- Byzantine-robust aggregation family in simple language}
\label{sec:robust-agg}

All of these methods try to answer the same question: ``out of the updates the
server got this round, which should I actually use?''

\textbf{\fedavg.}
Average everyone's updates, weighted by how much data each client had. No
defence. The baseline.

\textbf{\krum.}
Pick the single update that is closest to the other updates in distance.
Intuition: ``believe the one who is most like the crowd.'' Sensitive: if the
crowd drifts a little bit, Krum picks a bad update.

\textbf{\mkrum.}
Same scoring as Krum but average multiple ``most-like-the-crowd'' updates
instead of picking just one. More stable than Krum.

\textbf{\median (coordinate-wise).}
For each parameter of the model, take the median value across clients. A few
extreme values can't bias the median. Statistically principled but penalises
honest non-IID clients with legitimately unusual values.

\textbf{\tmean (coordinate-wise).}
For each parameter, drop the top and bottom 20\% of values and average the
rest. Similar to median in spirit; slightly more flexible.

\textbf{\bulyan.}
A two-stage defence: first filter outliers using a \mkrum-style step, then
average what's left using \tmean. Considered the strongest classical defence.

\textbf{The common limitation.}
All five treat each round independently. They know nothing about the client's
past. They are all geometry-based. None of them can catch a sleeper before it
activates, because the sleeper looks normal in every single round's geometry.

\textbf{What PoC adds.}
PoC is the first ``behavioural'' rule in this comparison grid. It doesn't look
at weight geometry; it looks at accuracy history. This is why it is the only
method that stays safe under the sleeper attack.

% ----------------------------------------------------------------------------
\subsection{Part B-VIII --- Backdoors, sleeper attacks, and ``critical-layer flip'' (plain English)}
\label{sec:backdoor-conceptual}

\textbf{What a backdoor is (not a classifier bug).}
A backdoor is a hidden rule inserted into the model: ``on most inputs, behave
normally, but on these specific inputs, output the attacker's chosen label.''
So the model still looks good on ordinary evaluation, while failing silently
on the attacker's target pattern.

\textbf{Sleeper behaviour, step by step.}
\begin{enumerate}
\item Rounds 1--14: the Byzantine client trains normally. Its updates look
like any other honest update.
\item Round 15: the client switches mode. It flips labels on the target class
(APB becomes LBBB in this project).
\item Rounds 15--40: it keeps uploading poisoned updates that nudge the global
model toward the backdoor.
\end{enumerate}
This is called a \emph{sleeper} because it is dormant for a long time before
activating.

\textbf{``Backdoor-critical layer'' in simple terms.}
The deeper layers of a classifier near the output act like a decision-maker;
the earlier layers extract features. A backdoor can be inserted by changing
only a small number of layers near the output (the ``critical'' ones). Result:
most of the model's weights look identical to the honest model, so geometric
defences see only a tiny difference, well within honest non-IID noise.

\textbf{Why this breaks robust aggregation.}
Robust aggregators look for updates that are ``too far'' from the crowd. A
critical-layer flip changes almost nothing in most coordinates, and only
slightly in a few. The update \emph{is} close to the crowd. Nothing to filter.

\textbf{Why PoC catches it.}
The attacker's own validation accuracy on honest data drops as soon as the flip
begins. That drop is exactly the signal PoC's EMA watches. Geometry sees no
anomaly; behaviour does.

% ----------------------------------------------------------------------------
\subsection{Part B-IX --- Classification metrics in simple language}
\label{sec:metrics-conceptual}

\textbf{Precision, recall, F1 in one picture.}
Suppose the model predicted some beats as APB.
\begin{itemize}
\item \textbf{Precision:} of the ones the model called APB, how many actually are APB?
\item \textbf{Recall:} of the beats that actually are APB, how many did the model catch?
\item \textbf{F1:} a single number that balances precision and recall (higher is better; $F1 = 2 \cdot \frac{P \cdot R}{P + R}$).
\end{itemize}

\textbf{Macro-F1.}
Average F1 across all classes, treating each class equally (ignoring how many
examples each class has). This is why it can look good even if the rare class
is failing.

\textbf{APB-F1 (our primary metric).}
The F1 score specifically on the APB class. Why primary? Because APB is rare
and clinically dangerous; missing APB is a safety failure. A method that gets
high macro-F1 but low APB-F1 is clinically useless.

\textbf{Backdoor Success Rate (BSR).}
Fraction of trigger samples that the classifier assigns to the attacker's
chosen wrong label. Lower is better. BSR matters in the diffusion-augmented
setting: two defences can have similar APB-F1 but very different BSR because
a higher BSR means the attacker has still leaked backdoor capacity even if the
average metric looks healthy.

\textbf{Why precision often matters more than recall here.}
In cardiology, a false APB triggers a workup that costs time and money; a
missed APB risks patient harm. Both matter, but because the dataset has so
many Normal beats, the model easily gets high recall by being liberal, and the
harder part is keeping precision up without sacrificing recall.

\textbf{One viva-ready line.}
\textit{``Macro-F1 hides rare-class failure; we use APB-F1 as the primary
clinical metric and BSR as the security metric because together they show both
the clinical outcome and the residual attack success.''}

% ----------------------------------------------------------------------------
\subsection{Part C --- Every design choice explained (``why this, not that'')}
\subsection{Part C --- Every design choice explained (``why this, not that'')}

\subsubsection{C.1 --- Why CNN--LSTM for the classifier?}

\textbf{What it does.} A 1D CNN extracts local morphology features (the
shape of the QRS complex, the position and amplitude of P and T waves).
An LSTM on top of the CNN features captures the temporal dependency
between samples of the 360-sample window.

\textbf{Why not alternatives?}
\begin{itemize}
\item \textbf{Pure MLP.} Ignores temporal locality; can work but is brittle on noisy waveforms and is not standard for ECG in the literature.
\item \textbf{Pure CNN.} Captures local morphology well, but a single-window ECG beat depends on the preceding context; LSTM adds temporal robustness for noisy or borderline beats.
\item \textbf{Pure RNN / GRU.} Similar to LSTM but typically slightly less stable on long morphology windows.
\item \textbf{Transformers.} Would be more flexible, but:
  \begin{itemize}
  \item They need more data per site (FL clients often have small splits).
  \item They have more hyperparameters, which would confound the security study.
  \item Compute cost during 40 rounds $\times$ 10 clients $\times$ 3 epochs is non-trivial for a research prototype on a single T4.
  \end{itemize}
\item \textbf{Pretrained foundation models.} Possible but they hide capacity: in an FL security study, you need the classifier to be simple enough that the security mechanism, not the base model, explains the results.
\end{itemize}

\textbf{The deeper reason.} The paper's contribution is a security /
governance mechanism. A classifier that is too powerful would \emph{mask}
the attacks (and reviewers would ask ``did your defence matter, or did
the classifier just absorb the attack?''). CNN--LSTM is the right level of
complexity: strong enough to give meaningful APB-F1, simple enough to
expose failure under attack.

\subsubsection{C.2 --- Why 1D latent diffusion for augmentation?}

\textbf{What it does.} A class-conditional 1D UNet with latent diffusion
generates synthetic ECG windows of a specific class on request.
Inference uses 30 DDPM steps. Conditioning is by channel concatenation:
the normalised class label is broadcast along the time axis and fed as
a second input channel.

\textbf{Why diffusion (vs GAN / VAE / TimeGAN / copy-paste oversampling)?}
\begin{itemize}
\item \textbf{GAN / TimeGAN.} The proposal originally considered TimeGAN. GANs have known training instability (mode collapse, vanishing gradients). For a federated setting, where training budgets are already tight, diffusion is empirically more stable and gives better class diversity.
\item \textbf{VAE.} Simpler to train but produces blurrier samples; class-conditional fidelity on rare arrhythmia morphology suffers.
\item \textbf{SMOTE / copy-paste resampling.} Does not produce genuinely new morphology; does not help the model learn richer APB variants; and is easy to exploit (an attacker can game oversampling trivially).
\item \textbf{Latent diffusion.} Compresses input to a lower-dimensional latent space, then diffuses there. This is much cheaper than pixel-space (sample-space) diffusion and is the same rationale as the 2D image case.
\end{itemize}

\textbf{Why 1D and not 2D?}
\begin{itemize}
\item ECG is natively a 1D signal; turning it into a 2D spectrogram and then back again adds artifacts and complexity without benefit.
\item 1D UNet is lighter, faster, and matches inductive bias.
\end{itemize}

\textbf{Why 30 DDPM steps?}
Fewer steps $\Rightarrow$ faster inference but worse quality;
more steps $\Rightarrow$ slower but diminishing returns.
30 steps was the lowest count at which generated waveforms still had
clinically plausible QRS morphology, chosen empirically.

\textbf{Why 500 synthetic samples per request?}
Large enough to shift the loss landscape on the deficient class
during a training round, but small enough to avoid dominating the
local training set. A request cap prevents a malicious requester
from overwhelming honest data even if somehow approved.

\subsubsection{C.3 --- Why blockchain at all? And specifically why Ganache?}

\textbf{What problem the blockchain solves.}
In a multi-hospital deployment, \emph{who trusts whom} is the core
question. A simple ``signed append-only log'' maintained by one party
(say, the orchestrator) still requires participants to trust that party
to not rewrite history. Blockchain gives:
\begin{enumerate}
\item A shared, ordered, tamper-evident record of events.
\item Verifiable approval receipts (so the augmentation pipeline becomes auditable).
\item A neutral substrate on which post-hoc forensics can run, without any single participant being able to rewrite the past.
\end{enumerate}

In this project the ledger has exactly two roles:
\begin{enumerate}
\item Persist update and round-completion receipts (what was aggregated and when).
\item Persist the synthetic-request approval chain (what was approved, by whom, at what PoC score).
\end{enumerate}

\textbf{Why Ganache (local Ethereum) for a research prototype?}
\begin{itemize}
\item \textbf{Deterministic and free.} Every run is reproducible; no gas cost limits experimentation.
\item \textbf{Same EVM semantics as mainnet.} Solidity contracts, event logs, transaction receipts all behave like ``real'' Ethereum.
\item \textbf{Latency is predictable.} This keeps the federated training loop honest: if something takes time, it is the training, not a public mempool.
\item \textbf{No external dependency during experiments.} Stability matters for a 40-round experiment.
\end{itemize}

\textbf{Why an event-emitting (``LOG opcode'') contract, not state-heavy storage?}
\begin{itemize}
\item On Ethereum, persistent storage (\texttt{SSTORE}) is expensive; \texttt{LOG} opcodes (events) are cheap.
\item The receipts only need to be \emph{readable} to a validator who replays the chain; they don't need to be queryable via contract state.
\item Therefore the contract emits \texttt{UpdateLogged}, \texttt{RoundCompleted}, \texttt{SyntheticRequested/Approved/Generated} events, and keeps almost no persistent state.
\item This is what enables the 84.6\% gas reduction figure vs storing arrays via SSTORE.
\end{itemize}

\textbf{Why not Hyperledger / permissioned blockchain?}
\begin{itemize}
\item Hyperledger Fabric would actually fit a hospital consortium very well in production.
\item For a research prototype, Ganache was chosen because it is faster to set up, deterministic, and aligned with the widely-understood EVM model, which makes the contribution reproducible.
\item The paper explicitly calls permissioned chains / L2 a reasonable production direction.
\end{itemize}

\textbf{Key boundary (do not over-claim).}
The blockchain is \emph{not} performing training; PoC is \emph{not} a
consensus protocol; the chain does not replace the orchestrator. It is
an accountability layer.

\subsubsection{C.4 --- Why Proof-of-Contribution (behavioural) and not robust aggregation (geometric)?}

\textbf{Core insight.} Geometry-only defences answer the question
``does this client's update look like the crowd in this round?''
This is blind to \emph{time}.

A sleeper attacker by design looks exactly like the crowd until round
15, then silently flips. Up to round 15, no per-round rule can filter
it --- there is no anomaly yet. After round 15, the poisoned updates
can still be crafted to stay close to honest updates in weight space,
especially under non-IID where honest clients legitimately disagree.

What geometry cannot see but PoC can:
\begin{itemize}
\item \textbf{Time-varying behaviour.} PoC's EMA is built precisely to detect a change in behaviour over multiple rounds.
\item \textbf{Self-inflicted accuracy drop.} A label-flip attacker hurts its own validation accuracy on honest validation data. Geometry doesn't know that; PoC is exactly that signal.
\end{itemize}

\textbf{Why not a heavier reputation system (like Kang et al.)?}
\begin{itemize}
\item Those systems are designed for participation economics / identity management. They are heavier and come with their own attack surface.
\item PoC is intentionally minimal (one EMA, one participation factor) because the aim was to prove the \emph{principle}: a behavioural history is enough.
\end{itemize}

\subsubsection{C.5 --- Why the same model family across all experiments?}

This is an important viva question. The honest answer:
\begin{itemize}
\item \textbf{Experimental control.} To argue ``defence X is better than defence Y'', everything else must be identical: same classifier, same data partitioning, same attacks, same non-IID setting, same training budget. Otherwise the conclusion is contaminated.
\item \textbf{Fairness to baselines.} Changing the classifier for each defence would let someone claim your gains come from model capacity, not from the defence.
\item \textbf{Clarity of the contribution.} The research contribution is the security/governance stack, not the classifier or the generator.
\end{itemize}

If asked ``did you try different classifiers / different generators?'':
answer honestly that controlled comparison is why a single well-chosen
pair (CNN--LSTM + 1D latent UNet) was held constant; exploring other
classifiers or generators is legitimate future work.

\subsubsection{C.6 --- Why these specific baselines?}

\begin{itemize}
\item \textbf{FedAvg.} The ``no-defence'' reference. Any honest security claim must include it to show what happens without defence.
\item \textbf{Krum.} The canonical single-update Byzantine-robust aggregator, known to be geometry-based; expected to fail under sleeper.
\item \textbf{Multi-Krum.} Generalisation of Krum to multi-client averaging; less brittle than Krum but still geometric.
\item \textbf{Coordinate-wise Median.} Statistically optimal under certain Byzantine ratios; penalises honest non-IID tails.
\item \textbf{Trimmed-Mean.} Closely related; drops extreme values per coordinate; penalises non-IID outliers.
\item \textbf{Bulyan.} Composes Multi-Krum with Trimmed-Mean; strongest classical robust aggregator.
\item \textbf{PoC (ours).} Behavioural.
\end{itemize}

Together they span the three natural families (selection-based,
coordinate-wise, composed) and make it impossible for a reviewer to
say ``you didn't test against the obvious baseline''.

\subsubsection{C.7 --- Why these specific attacks?}

The three attacks are deliberately of \emph{three different flavours},
not three variants of the same thing:
\begin{enumerate}
\item \textbf{Gaussian poisoning.} The blatant ``sanity-check'' attack from the original robust-FL literature. If a defence cannot stop this, it's unusable. FedAvg famously fails here.
\item \textbf{Static label-flip (Normal$\leftrightarrow$LBBB, colluding).} A stealthier attack where updates \emph{look} like valid training steps (normal magnitudes, normal directions) but on wrong labels. This specifically probes whether a defence relies only on weight-space geometry.
\item \textbf{Sleeper backdoor (APB$\rightarrow$LBBB, critical-layer fine-tune, round 15 activation).} The worst case for any per-round defence: pre-activation there is no anomaly, so geometry cannot pre-train a filter; post-activation only a subset of backdoor-critical layers is modified, keeping weight-space signatures normal.
\end{enumerate}

This triple is also why the paper can claim ``across all three attacks''
with a straight face: they span the adversary's strategy space.

\subsubsection{C.8 --- Why the specific hyperparameters?}

\begin{itemize}
\item \textbf{Dirichlet $\alpha=0.5$.} Standard non-IID stress setting in FL; not so extreme that training collapses even for honest FedAvg, not so mild that it's essentially IID.
\item \textbf{10 clients, 2 Byzantine (20\%).} Realistic minority-threat ratio; matches classical robust-FL assumptions (Krum-family is designed for <25\% Byzantines).
\item \textbf{40 rounds.} Enough for global convergence and also enough headroom for the sleeper to activate at round 15 and have 25 more rounds to manifest. Fewer rounds would hide the sleeper's impact.
\item \textbf{3 local epochs, batch 128, Adam $10^{-3}$.} Conservative, stable training.
\item \textbf{EMA $\gamma=0.7$.} Chosen so that a single bad round can move the EMA noticeably, but a single lucky round cannot rehabilitate a bad client.
\item \textbf{Trust threshold $\tau=0.4$.} Empirically the level at which pre-activation attackers are admitted (so they don't poison the very first rounds via denial of synthetic help) but post-activation attackers are rejected.
\item \textbf{Top-$K=7$ out of 10.} Enough to keep honest non-IID diversity, but still filters out the bottom three lowest-PoC clients per round.
\end{itemize}

\subsubsection{C.9 --- Why these metrics?}

\begin{itemize}
\item \textbf{Primary: APB-F1.} Clinical relevance. APB is rare and dangerous; missing APB is a safety failure, not a performance issue.
\item \textbf{Secondary: Macro-F1.} Required by tradition, but deliberately de-emphasised because it can mask minority-class failure.
\item \textbf{Backdoor Success Rate (BSR).} Only applies in the diffusion-augmented setting. It captures residual backdoor capacity even when APB-F1 is healthy. This is where the blind-vs-trust-gated difference shows up most clearly.
\item \textbf{Latency.} Included only to show the ledger is not a wall-clock bottleneck; explicitly not a contribution.
\end{itemize}

% ----------------------------------------------------------------------------
\subsection{Part D --- Threat model, stated formally (for questions)}

\textbf{Adversary capabilities.}
\begin{itemize}
\item Controls 2 out of 10 clients throughout all 40 rounds (20\% Byzantine, static).
\item Knows the training protocol, defences, and the existence of PoC.
\item Can collude between its 2 clients (synchronise batch truncation, synchronise activation).
\item Can choose between three attack strategies (Gaussian, label-flip, sleeper) and, in the augmented setting, can also file flipped-label synthetic requests.
\end{itemize}

\textbf{Adversary limitations.}
\begin{itemize}
\item Cannot corrupt the orchestrator.
\item Cannot rewrite chain history.
\item Cannot tamper with the public proxy validation set used for PoC accuracy scoring.
\item Does not adaptively evolve PoC-aware attacks --- this is explicitly listed as a limitation and future work.
\end{itemize}

% ----------------------------------------------------------------------------
\subsection{Part E --- The result story, retold as a narrative}

Think of the results as answering five questions in sequence:

\begin{enumerate}
\item \emph{Is any standard FL safe here?} No. FedAvg collapses to APB-F1 $=0$ under Gaussian.
\item \emph{Can single-round geometry save us?} Krum collapses to 0.306 under sleeper. Multi-Krum, Median, Trimmed-Mean, Bulyan all drop below 0.75 in at least one attack.
\item \emph{Does behavioural history help?} Yes. PoC is the \emph{only} defence that stays above 0.75 across all three attacks (worst case 0.786).
\item \emph{Is there a second attack surface nobody looks at?} Yes --- the augmentation pipeline. An attacker can inject poison via flipped-label synthetic requests.
\item \emph{Can governance close it?} Yes. Trust-gated diffusion with PoC raises APB-F1 to 0.879 and lowers the final-round backdoor success rate vs blind diffusion.
\end{enumerate}

This sequence is exactly what you should walk a jury through.

% ----------------------------------------------------------------------------
\subsection{Part F --- Mega Q\&A (more than 50)}

\textit{If you can answer these, you are in hero mode.}

\subsubsection{F.1 --- Model choices}
\begin{enumerate}
\item \textbf{Why CNN--LSTM and not Transformer?} Data efficiency, stability, prototype clarity, reduced confounding.
\item \textbf{Why not use a pretrained ECG foundation model?} It would mask the security effects the paper is trying to isolate.
\item \textbf{Why use LSTM at all; isn't CNN enough?} LSTM adds temporal robustness on borderline windows.
\item \textbf{How many parameters is the classifier?} A compact CNN--LSTM in the low millions; exact number is irrelevant to the security claim; the point is it's small enough for all 10 FL clients to train.
\item \textbf{Why 360-sample windows at 360Hz?} MIT-BIH is at 360Hz and EC57 recommends single-beat windows; 360 samples = 1 second ~centred on the R-peak.
\end{enumerate}

\subsubsection{F.2 --- Diffusion choices}
\begin{enumerate}
\item \textbf{Why diffusion, not a GAN?} Diffusion is more stable and less prone to mode collapse; matters because the generator is trained once and trusted later.
\item \textbf{Why latent diffusion vs pixel-space?} Cheaper inference, same empirical quality in 1D.
\item \textbf{Why 30 steps?} Minimum step count at which QRS morphology looked clinically plausible.
\item \textbf{Is the generator trained during FL?} No. It is pretrained once on pooled honest ECG data and then frozen; only inference runs in the FL loop.
\item \textbf{Does the generator see test data?} No; only pooled honest training data.
\item \textbf{Why class-conditional by channel concatenation?} Simplest robust conditioning method for a 1D UNet; avoids embedding networks.
\item \textbf{Why cap at 500 synthetic samples per request?} Enough to repair rare-class deficit; not so many that honest data is overwhelmed.
\item \textbf{What stops an attacker from asking too many times?} On-chain rate-limiting via the same PoC gate; rejected requests are logged.
\end{enumerate}

\subsubsection{F.3 --- FL choices}
\begin{enumerate}
\item \textbf{Why 10 clients?} Enough diversity for meaningful non-IID; feasible for 40 rounds $\times$ 3 epochs on a single T4 with synchronous diffusion inference.
\item \textbf{Why 40 rounds?} Long enough for convergence and post-sleeper dynamics to show up.
\item \textbf{Why 3 local epochs?} Classic FL compromise --- fewer epochs mean more communication, more epochs risk client drift.
\item \textbf{Why Adam and not SGD?} Stable on ECG with limited tuning; removes a free variable from the security experiment.
\item \textbf{Why top-$K=7$?} Filters out a modest bottom tail without over-pruning the honest non-IID diversity.
\item \textbf{Does the server train anything?} No. The server only aggregates and evaluates.
\item \textbf{Is PoC computed during training or after?} It is computed each round, before selection for the next round.
\end{enumerate}

\subsubsection{F.4 --- Blockchain choices}
\begin{enumerate}
\item \textbf{Why blockchain at all?} Tamper-evident, shared receipts; auditable approvals.
\item \textbf{Why not a signed log?} A signed log still needs one party to maintain it.
\item \textbf{Why Ganache?} Deterministic local Ethereum; same EVM semantics as production.
\item \textbf{Why event-only contract?} Events are cheap; receipts don't need contract-level queryability.
\item \textbf{Does the chain slow training?} No; overhead is bounded; latency is not claimed as a contribution.
\item \textbf{Is the chain needed for correctness?} For correctness of FL, no. For \emph{auditability} and the augmentation-governance property, yes.
\item \textbf{Could you swap Ganache for Hyperledger?} Yes; permissioned chains fit a real hospital consortium.
\item \textbf{Is gas cost a real claim?} No. It is only a sanity check that the chosen design (events, not storage) would not cripple a production deployment.
\item \textbf{What if the chain goes down?} Defensive handlers treat that round as neutral and training continues.
\end{enumerate}

\subsubsection{F.5 --- Security and attacks}
\begin{enumerate}
\item \textbf{What is the threat model?} 20\% colluding Byzantine clients, static throughout training, with full knowledge of protocol.
\item \textbf{What is a sleeper attack?} A client that acts honestly until the global model stabilises, then flips labels.
\item \textbf{What is ``backdoor-critical-layer flip''?} Only the tail layers of the classifier are fine-tuned on the flipped labels, so most weights look unchanged.
\item \textbf{What is the augmentation stealth channel?} A malicious synthetic request with a flipped label.
\item \textbf{Why can't Krum block the augmentation channel?} Because the poison is baked into training data before the aggregator ever sees any update.
\item \textbf{How does PoC catch the sleeper?} Flipping labels lowers the attacker's own validation accuracy; PoC's EMA reacts and drops the attacker out of top-$K$ and below $\tau$.
\item \textbf{What happens between rounds 15--25?} APB-F1 temporarily oscillates for some defences; for PoC it recovers; for Krum it degrades.
\item \textbf{What is BSR?} Backdoor Success Rate: fraction of triggered samples classified as the attacker-desired label.
\item \textbf{Why is BSR high even when APB-F1 is high?} Because APB-F1 is a whole-class metric; BSR catches targeted failures that F1 can average away.
\end{enumerate}

\subsubsection{F.6 --- Metrics, experiments, honesty}
\begin{enumerate}
\item \textbf{Why APB-F1 as primary?} Clinical relevance; minority-class safety.
\item \textbf{Why macro-F1 as secondary?} Tradition and sanity; shows that the per-class story is the real story.
\item \textbf{Why single-seed?} Compute cost of 40-round FL with synchronous latent-diffusion inference; multi-seed is future work; the observed effect is structural.
\item \textbf{Is there variance analysis?} No multi-seed variance; explicitly declared as a limitation.
\item \textbf{Do you cross-validate?} Each client holds out 15\% for test; test is held out across all 40 rounds; no snooping.
\item \textbf{What about PTB-XL / MIMIC-IV?} Not evaluated here; future work.
\item \textbf{Are the numbers in prose and figures consistent?} Yes; they are extracted programmatically from the serialised artefacts; this addresses an ACISP reviewer complaint.
\end{enumerate}

\subsubsection{F.7 --- Limitations and honest caveats}
\begin{enumerate}
\item Single-seed per cell (compute budget).
\item Static 20\% Byzantine ratio; higher ratios not tested.
\item Non-adaptive adversary (does not evolve with PoC knowledge).
\item Only MIT-BIH; no cross-dataset validation.
\item Gas accounting is illustrative, not a production benchmark.
\item No formal proof that PoC is optimal; it is a strong empirical baseline for a behavioural signal.
\item Orchestrator is assumed honest; decentralising the orchestrator is future work.
\item Public proxy validation set could itself be attacked in principle; production would need signed proxy or secure evaluation.
\end{enumerate}

% ----------------------------------------------------------------------------
\subsection{Part G --- Common viva traps and how to avoid them}

\begin{itemize}
\item \textbf{Trap:} ``So the blockchain does the learning / consensus?''\\
\textbf{You:} No. The ledger records receipts. PoC is a score, not a proof. Training is off-chain.

\item \textbf{Trap:} ``Macro-F1 says FedAvg is fine under label-flip.''\\
\textbf{You:} True in aggregate, but macro-F1 hides the minority-class failure; the whole point is that APB-F1 is the deployment-critical number.

\item \textbf{Trap:} ``Krum is supposed to be robust; why did it fail?''\\
\textbf{You:} Krum is robust to geometric outliers in a single round. Sleeper attacks are time-varying and stealthy; they don't look like outliers pre-activation, and post-activation the critical-layer flips keep the weights geometrically normal.

\item \textbf{Trap:} ``Isn't diffusion overkill for ECG?''\\
\textbf{You:} It is heavier than SMOTE, but diffusion gives genuine class-conditional diversity that SMOTE or GAN cannot reliably produce; and more importantly, the paper's novelty is not the generator but governance around the generator.

\item \textbf{Trap:} ``What if attackers compromise the orchestrator?''\\
\textbf{You:} That is outside the threat model, but the ledger would still expose the manipulation post-hoc; decentralising the orchestrator is future work.

\item \textbf{Trap:} ``Why didn't you test with more attackers?''\\
\textbf{You:} 20\% is the standard setting for robust-FL comparisons. Higher ratios (explicit limitation) are future work.

\item \textbf{Trap:} ``Why not just use a bigger classifier to beat the attack?''\\
\textbf{You:} Capacity doesn't help against label-flip / backdoor attacks; in fact, higher-capacity models memorise backdoors more reliably.
\end{itemize}

% ----------------------------------------------------------------------------
\subsection{Part H --- 60-second elevator pitch (memorise verbatim)}

\textit{``Clinical ECG federated learning has two simultaneous problems: 
non-IID hospital data and realistic Byzantine clients that target rare 
dangerous classes. Classical robust aggregation only looks at one 
round of weights at a time, which is why sleeper backdoor attacks 
defeat them. Active-Ledger fixes this with a behavioural trust score 
(Proof-of-Contribution) that uses multi-round accuracy history, 
stored on a tamper-evident ledger. The same trust score also 
governs a separate attack surface nobody else has closed --- 
malicious synthetic data requests. Across three different attacks, 
Proof-of-Contribution is the only defence that keeps APB-F1 above 
the 0.75 deployability line, and adding trust-gated diffusion raises 
it further to 0.879 while lowering backdoor success rate. The 
blockchain is explicitly not doing the learning --- it is the audit 
layer that makes both aggregation and augmentation accountable.''}

% ----------------------------------------------------------------------------
\subsection{Part I --- What to study the night before (priority list)}

\begin{enumerate}
\item The 8-step round trace (Part B) --- can you say all 8 steps out loud?
\item The five-question result narrative (Part E).
\item The three attack descriptions, each in one sentence.
\item Why geometry-only fails against sleeper (Part C.4).
\item Why the ledger is needed (Part C.3) --- without over-claiming.
\item The six memorised numbers: 0.75, 0.786, 0.879, 0.857, 0.306, 0.000.
\item The elevator pitch (Part H).
\item The top 20 Q\&A answers from Section 9.
\end{enumerate}

\textit{If you have all of this, you can defend every slide.}

\end{document}

